\subsection{Procedimiento de clasificación y presentación de resultados}

Una vez aplicado el instrumento de la Tabla~\ref{tab:kano_preguntas} a las cinco personas reales del dominio, cada pareja de respuestas se codifica con la Tabla~\ref{tab:kano_matriz}. La categoría final de cada característica $K1$ a $K8$ se obtiene por mayoría simple; si hay empate, se aplica el orden de prioridad M, O, A, I, R y Q. El resultado se registra en \texttt{datos/kano\_respuestas.csv} y se resume en una tabla con las columnas: característica, categoría asignada y frecuencias de M, O, A, I, R y Q.

\begin{table}[ht]
\centering
\footnotesize
\caption{Estructura de presentación de los resultados Kano (sin datos reales aún).}
\label{tab:resultado_kano}
\begin{tabular}{@{}llc@{}}
\toprule
Característica & Categoría final & Frecuencias (M/O/A/I/R/Q) \\
\midrule
K1 & \emph{Pendiente de datos reales} & --- \\
K2 & \emph{Pendiente de datos reales} & --- \\
K3 & \emph{Pendiente de datos reales} & --- \\
K4 & \emph{Pendiente de datos reales} & --- \\
K5 & \emph{Pendiente de datos reales} & --- \\
K6 & \emph{Pendiente de datos reales} & --- \\
K7 & \emph{Pendiente de datos reales} & --- \\
K8 & \emph{Pendiente de datos reales} & --- \\
\bottomrule
\end{tabular}
\end{table}

No se reporta ninguna categoría definitiva mientras no existan las respuestas reales de los cinco participantes. Cualquier clasificación presentada antes de ese paso sería evidencia fabricada y se excluye de este informe.
